\section{Analytic completions and the coupled certificate gauge}
\label{sec:completion-gauge}

The physical datum is prescribed only on the primitive lattice, including
zeros at absent primitive rays.  Values at nonprimitive coordinates do not
change the physical sum.  They can nevertheless change every term of a
triangle-inequality certificate.

\subsection{Declared finite completions}

Fix \(U,V\in\N\) and a rectangle
\[
 B_{U,V}=[1,U]\times[1,V]
\]
containing the primitive physical datum
\[
 k:\cP_2\cap B_{U,V}\longrightarrow\mathbb C.
\]

\begin{definition}[Completion class]
\label{def:completion-class}
Let \(\Ext_{U,V}(k)\) be the affine space of kernels
\(F:\N^2\to\mathbb C\) such that
\[
 \supp F\subseteq B_{U,V},
 \qquad
 F(u,v)=k(u,v)
 \quad\bigl((u,v)\in\cP_2\cap B_{U,V}\bigr).
\]
The values at nonprimitive sites are free analytic variables.  The
support-faithful zero completion
\[
 K^0(u,v)=
 \begin{cases}
  k(u,v),&(u,v)\in\cP_2\cap B_{U,V},\\
  0,&\text{otherwise}
 \end{cases}
\]
belongs to this class.
\end{definition}

Define the primitive signed functional
\begin{equation}
 \cS_{\rm prim}(k)
 =
 \sum_{\substack{u\le U,\ v\le V\\(u,v)=1}}
 \mu(u)\mu(v)k(u,v).
 \label{eq:primitive-functional}
\end{equation}

\begin{theorem}[Completion invariance]
\label{thm:completion-invariance}
For every \(F\in\Ext_{U,V}(k)\),
\begin{equation}
 \boxed{
 \cS_{\rm prim}(k)
 =
 \sum_{\substack{u,v\ge1\\(u,v)=1}}
 \mu(u)\mu(v)F(u,v).}
 \label{eq:completion-invariance}
\end{equation}
Thus the physical functional is constant on the complete affine class
\(\Ext_{U,V}(k)\).
\end{theorem}

\begin{proof}
The summation in \eqref{eq:completion-invariance} samples only primitive
coordinates, where every member of the class equals \(k\).
\end{proof}

\subsection{The full completion cost}

For \(d\ge1\), put
\[
 M_d(x)=\sum_{\substack{n\le x\\(n,d)=1}}\mu(n),
 \qquad
 \mathfrak M_d(T)
 =\max_{1\le n\le\lfloor T\rfloor}|M_d(n)|,
 \qquad
 F^{[d]}(x,y)=F(dx,dy).
\]
We set \(\mathfrak M_d(T)=0\) for \(T<1\).
For \(D\ge1\), define the exact weighted completion cost
\begin{align}
 \cC_D(F)
 :={}&
 \sum_{d\le D}|\mu(d)|
 \sum_{x,y\ge1}
 |M_d(x)M_d(y)|
 |\Delta_{12}F^{[d]}(x,y)|
 \notag\\
 &+
 \cT_D(F),
 \label{eq:full-completion-cost}
\end{align}
where \(\cT_D\) is the exact tail in
\eqref{eq:diagonal-tail}.

\begin{theorem}[Completion-gauge certificate]
\label{thm:completion-gauge}
For every declared completion class and every \(D\ge1\),
\begin{equation}
 \boxed{
 |\cS_{\rm prim}(k)|
 \le
 \inf_{F\in\Ext_{U,V}(k)}\cC_D(F).}
 \label{eq:completion-gauge}
\end{equation}
For the canonical zero completion,
\begin{equation}
 \cC_D(K^0)=\cC_\square(K^0)
 \qquad(D\ge1).
 \label{eq:canonical-cost-collapse}
\end{equation}
\end{theorem}

\begin{proof}
Fix \(F\in\Ext_{U,V}(k)\).  By
\cref{thm:completion-invariance}, its primitive sum is the physical
functional.  The exact TPC--71 diagonal divisor lift
\citep{WangTPC71} followed by
two-dimensional Abel summation gives
\[
 \cS_{\rm prim}(k)
 =
 \sum_{d\le D}\mu(d)
 \sum_{x,y\ge1}
 M_d(x)M_d(y)\Delta_{12}F^{[d]}(x,y)
 +\cR_D(F),
 \qquad
 |\cR_D(F)|\le\cT_D(F).
\]
Taking absolute values gives
\[
 |\cS_{\rm prim}(k)|\le\cC_D(F).
\]
Take the infimum.  For \(K^0\), all \(d\ge2\) slices and the tail vanish by
\cref{thm:primitive-support-collapse}; the \(d=1\) term is exactly
\eqref{eq:weighted-plaquette-cost}.
\end{proof}

\begin{remark}[Infimum versus a physical construction]
\label{rem:completion-quantifier}
The infimum is an exact finite-dimensional analytic benchmark.  To use a
completion in a uniform physical theorem, one must construct it before any
later extremizing coefficient vector is selected and control its cost
uniformly in all literal labels.  An optimizer chosen after seeing the
operator extremizer changes the quantifier.  Moreover, nonprimitive values
are analytic fill-ins and are never added to the physical reference mass.
\end{remark}

\begin{remark}[No separate monotonicity]
The \(d=1\) variation and the higher-diagonal tail can move in opposite
directions when a completion changes.  Neither term is an invariant of
the primitive datum.  Only the original physical functional and a proved
bound for the complete cost are invariant conclusions.
\end{remark}

\subsection{A coarse completion norm}

The weighted cost is the full location-sensitive deterministic majorant.
A coarser but sometimes convenient completion cost is
\begin{equation}
 \cE_D(F)
 :=
 \sum_{d\le D}|\mu(d)|
 \mathfrak M_d(U/d)\mathfrak M_d(V/d)
 \cV_\square(F^{[d]})
 +
 \cT_D(F).
 \label{eq:coarse-completion-cost}
\end{equation}
Then \(\cC_D(F)\le\cE_D(F)\).  The two displayed pieces in
\eqref{eq:coarse-completion-cost} must be evaluated for the same completion;
minimizing them independently would splice incompatible analytic objects.
